\documentclass[10pt,a4paper]{amsart}
\usepackage{amsmath,amssymb,amsthm,graphicx}
\usepackage{verbatim}
\newenvironment{reference}{\comment}{\endcomment}
%\newenvironment{reference}{}{}
\newenvironment{slogan}{\comment}{\endcomment}
\newenvironment{history}{\comment}{\endcomment}

% For commutative diagrams we use Xy-pic
\usepackage[all]{xy}

% We use 2cell for 2-commutative diagrams.
\xyoption{2cell}
% Theorem environments.
%
\theoremstyle{plain}
\newtheorem{theorem}[subsection]{Theorem}
\newtheorem{proposition}[subsection]{Proposition}
\newtheorem{lemma}[subsection]{Lemma}

\theoremstyle{definition}
\newtheorem{definition}[subsection]{Definition}
\newtheorem{example}[subsection]{Example}
\newtheorem{exercise}[subsection]{Exercise}
\newtheorem{situation}[subsection]{Situation}

\theoremstyle{remark}
\newtheorem{remark}[subsection]{Remark}
\newtheorem{remarks}[subsection]{Remarks}

\numberwithin{equation}{subsection}

% Macros
%
\def\lim{\mathop{\mathrm{lim}}\nolimits}
\def\colim{\mathop{\mathrm{colim}}\nolimits}
\def\Spec{\mathop{\mathrm{Spec}}}
\def\Hom{\mathop{\mathrm{Hom}}\nolimits}
\def\Ext{\mathop{\mathrm{Ext}}\nolimits}
\def\SheafHom{\mathop{\mathcal{H}\!\mathit{om}}\nolimits}
\def\SheafExt{\mathop{\mathcal{E}\!\mathit{xt}}\nolimits}
\def\Sch{\mathit{Sch}}
\def\Mor{\mathop{\mathrm{Mor}}\nolimits}
\def\Ob{\mathop{\mathrm{Ob}}\nolimits}
\def\Sh{\mathop{\mathit{Sh}}\nolimits}
\def\NL{\mathop{N\!L}\nolimits}
\def\CH{\mathop{\mathrm{CH}}\nolimits}
\def\proetale{{pro\text{-}\acute{e}tale}}
\def\etale{{\acute{e}tale}}
\def\QCoh{\mathit{QCoh}}
\def\Ker{\mathop{\mathrm{Ker}}}
\def\Im{\mathop{\mathrm{Im}}}
\def\Coker{\mathop{\mathrm{Coker}}}
\def\Coim{\mathop{\mathrm{Coim}}}

% Boxtimes
%
\DeclareMathSymbol{\boxtimes}{\mathbin}{AMSa}{"02}

%
% Macros for moduli stacks/spaces
%
\def\QCohstack{\mathcal{QC}\!\mathit{oh}}
\def\Cohstack{\mathcal{C}\!\mathit{oh}}
\def\Spacesstack{\mathcal{S}\!\mathit{paces}}
\def\Quotfunctor{\mathrm{Quot}}
\def\Hilbfunctor{\mathrm{Hilb}}
\def\Curvesstack{\mathcal{C}\!\mathit{urves}}
\def\Polarizedstack{\mathcal{P}\!\mathit{olarized}}
\def\Complexesstack{\mathcal{C}\!\mathit{omplexes}}
% \Pic is the operator that assigns to X its picard group, usage \Pic(X)
% \Picardstack_{X/B} denotes the Picard stack of X over B
% \Picardfunctor_{X/B} denotes the Picard functor of X over B
\def\Pic{\mathop{\mathrm{Pic}}\nolimits}
\def\Picardstack{\mathcal{P}\!\mathit{ic}}
\def\Picardfunctor{\mathrm{Pic}}
\def\Deformationcategory{\mathcal{D}\!\mathit{ef}}
\usepackage{cleveref}
\crefname{originalequation}{original equation}{original equations}
\Crefname{originalequation}{Original equation}{Original equations}
\makeatletter
\newlabel{equation-map-representable-into-presheaf}{{090F}{source}{}{sourcepointer-1}{}}
\newlabel{equation-map-representable-into-presheaf@cref}{{[originalequation][0][]090F}{[1][0][]source}{}{}{}}
\makeatother
\title{Injective and surjective maps of presheaves are the monomorphisms and epimorphisms}
\author{The Stacks Project Authors}
\date{stacks-project:sites.tex:local-snapshot-20261012; GNU FDL 1.2 or later}
\begin{document}
\maketitle
\noindent\textit{Editorial note.} All delivered mathematical text is the author's own, copied from the pinned Stacks Project chapter file sites-cohomology.tex; no statement or proof step was rewritten, reordered or removed.\par\medskip
\noindent\textit{Editorial note.} The proof is the author's complete argument; every step the author writes out is inside the counted span and no closing sentence is a placeholder.\par\medskip
\noindent\textit{Editorial note.} Every labelled reference occurring inside the counted spans is carried in assembly.external\_references as a non-load-bearing navigation pointer (same-chapter ones at their definition site, cross-chapter ones at the invocation inside the counted span), so no load-bearing mathematics is replaced by a stub. No bibliography citation occurs inside the counted spans.\par\medskip
\noindent\textit{Editorial note.} The checked chapter list for same-name collisions is the set of stacks chapters staged in this repository: chow, cohomology, divisors, etale-cohomology, flat, more-algebra, more-morphisms, morphisms, spaces-more-morphisms, varieties together with this card's chapter sites-cohomology.\par\medskip
\noindent\textit{Editorial note.} The preamble is the chapter's own retained minimal preamble (sources/stacks-sites-cohomology/retained-preamble.tex), a verbatim excerpt of the native preamble: the Xy-pic setup and the native theorem declarations with the mathematical macros.\par\medskip
\noindent\textit{Allowed standard prerequisites (editorial).}\begin{itemize}
\item Injective and surjective maps of presheaves. Same chapter, Definition definition-presheaves-injective-surjective, recorded as the context of this card. Supplies the two notions compared with monomorphisms and epimorphisms.
\item Representable presheaves and the Yoneda identification. Same chapter, the identification of the elements of a presheaf at an object with natural transformations from the representable presheaf, used at source lines 389-396. Turns elements of a value into natural transformations so that monicity can be applied.
\item Pushouts in the category of sets. Standard: the pushout of a map of sets along itself has two canonical maps from the target, which coincide exactly when the map is surjective. Produces the two natural transformations whose equality forces surjectivity.
\end{itemize}\medskip
\ifdefined\hypertarget\hypertarget{sourcepointer-1}{}\fi
\noindent\textit{Source reference.} Original equation 090F: sources/stacks-sites/pinned-original-source.tex, source lines 390--390. Navigation pointer invoked at proof line 390; the counted span contains the pointer only and the cited result is not reproduced inside it, so the pointer is registered as a non-load-bearing navigation stub in the established form for Stacks material. No bibliography citation is involved.\par\medskip
\begin{definition}
\label{definition-presheaves-injective-surjective}
Let $\mathcal{C}$ be a category, and let $\varphi : \mathcal{F}
\to \mathcal{G}$ be a map of presheaves of sets.
\begin{enumerate}
\item We say that $\varphi$ is {\it injective} if for every object
$U$ of $\mathcal{C}$ the map $\varphi_U : \mathcal{F}(U)
\to \mathcal{G}(U)$ is injective.
\item We say that $\varphi$ is {\it surjective} if for every object
$U$ of $\mathcal{C}$ the map $\varphi_U : \mathcal{F}(U)
\to \mathcal{G}(U)$ is surjective.
\end{enumerate}
\end{definition}

\begin{lemma}
\label{lemma-mono-epi}
The injective (resp.\ surjective) maps defined above
are exactly the monomorphisms (resp.\ epimorphisms) of
$\textit{PSh}(\mathcal{C})$. A map is an isomorphism
if and only if it is both injective and surjective.
\end{lemma}

\begin{proof}
We shall show that $\varphi : \mathcal{F} \to
\mathcal{G}$ is injective if and only if it is a monomorphism
of $\textit{PSh}(\mathcal{C})$. Indeed, the ``only if''
direction is straightforward, so let us show the ``if''
direction. Assume that $\varphi$ is a monomorphism. Let
$U \in \Ob(\mathcal{C})$; we need to show that $\varphi_U$ is
injective. So let $a, b \in \mathcal{F}(U)$ be such that
$\varphi_U (a) = \varphi_U (b)$; we need to check that $a = b$.
Under the isomorphism
(\ref{equation-map-representable-into-presheaf}), the elements
$a$ and $b$ of $\mathcal{F}(U)$ correspond to two natural
transformations
$a', b' \in \Mor_{\textit{PSh}(\mathcal{C})}(h_U, \mathcal{F})$.
Similarly, under the analogous isomorphism
$\Mor_{\textit{PSh}(\mathcal{C})}(h_U, \mathcal{G})
= \mathcal{G}(U)$,
the two equal elements $\varphi_U (a)$ and $\varphi_U (b)$ of
$\mathcal{G}(U)$ correspond to the two natural transformations
$\varphi \circ a', \varphi \circ b'
\in \Mor_{\textit{PSh}(\mathcal{C})}(h_U, \mathcal{G})$,
which therefore must also be equal. So
$\varphi \circ a' = \varphi \circ b'$, and thus $a' = b'$
(since $\varphi$ is monic), whence $a = b$. This finishes (1).

\medskip\noindent
We shall show that $\varphi : \mathcal{F} \to
\mathcal{G}$ is surjective if and only if it is an epimorphism
of $\textit{PSh}(\mathcal{C})$. Indeed, the ``only if''
direction is straightforward, so let us show the ``if''
direction. Assume that $\varphi$ is an epimorphism.

\medskip\noindent
For any two morphisms $f : A \to B$ and $g : A \to C$ in the
category $\textit{Sets}$, we let $\text{inl}_{f,g}$ and
$\text{inr}_{f,g}$ denote the two canonical maps from
$B$ and $C$ to $B \coprod_A C$. (Here, the pushout is
evaluated in $\textit{Sets}$.)

\medskip\noindent
Now, we define a presheaf $\mathcal{H}$ of sets on $\mathcal{C}$
by setting $\mathcal{H}(U)
= \mathcal{G}(U) \coprod_{\mathcal{F}(U)} \mathcal{G}(U)$ (where
the pushout is evaluated in $\textit{Sets}$ and induced by
the map $\varphi_U : \mathcal{F}(U) \to \mathcal{G}(U)$) for
every $U \in \Ob(\mathcal{C})$; its action on morphisms is
defined in the obvious way (by the functoriality of pushout).
Then, there are two natural
transformations $i_1 : \mathcal{G} \to \mathcal{H}$ and
$i_2 : \mathcal{G} \to \mathcal{H}$ whose components at an object
$U \in \Ob(\mathcal{C})$ are given by the maps
$\text{inl}_{\varphi_U, \varphi_U}$ and
$\text{inr}_{\varphi_U, \varphi_U}$, respectively. The
definition of a pushout shows that $i_1 \circ \varphi
= i_2 \circ \varphi$, whence $i_1 = i_2$ (since $\varphi$ is an
epimorphism). Thus, for every $U \in \Ob(\mathcal{C})$, we have
$\text{inl}_{\varphi_U, \varphi_U}
= \text{inr}_{\varphi_U, \varphi_U}$. Thus, $\varphi_U$
must be surjective (since a simple combinatorial argument shows
that if $f : A \to B$ is a morphism in $\textit{Sets}$, then
$\text{inl}_{f,f} = \text{inr}_{f,f}$ if and
only if $f$ is surjective). In other words, $\varphi$ is
surjective, and (2) is proven.

\medskip\noindent
We shall show that $\varphi : \mathcal{F} \to
\mathcal{G}$ is both injective and surjective if and only if it
is an isomorphism of $\textit{PSh}(\mathcal{C})$. This time,
the ``if'' direction is straightforward. To prove the ``only if''
direction, it suffices to observe that if $\varphi$ is both
injective and surjective, then $\varphi_U$ is an invertible map
for every $U \in \Ob(\mathcal{C})$, and the inverses of these
maps for all $U$ can be combined to a natural transformation
$\mathcal{G} \to \mathcal{F}$ which is an inverse to $\varphi$.
\end{proof}

\end{document}
